%%%%%%%%%%%%%%%%%%%%%%%%%%%%%%%%%%
% Kurzfassung als eigenstaendiges Dokument
% Layout identisch zur Kurzfassung in master.pdf (Seite vi)
% Praeambel-Einstellungen entsprechen den layoutrelevanten Teilen von config.tex
%%%%%%%%%%%%%%%%%%%%%%%%%%%%%%%%%%

\documentclass[
    fontsize=12pt,
    BCOR=5mm,
    DIV=12,
    parskip=half,
    paper=a4,
    numbers=noenddot
]{scrreprt}

\usepackage[utf8]{inputenc}
\usepackage[T1]{fontenc}
\usepackage[ngerman]{babel}
\usepackage{lmodern}
\usepackage{setspace}

% Schriften wie in config.tex
\renewcommand*{\familydefault}{\sfdefault}
\addtokomafont{disposition}{\sffamily}

% Kopf-/Fusszeile wie in config.tex
\RequirePackage[automark]{scrlayer-scrpage}
\renewcommand{\chaptermarkformat}{}
\RedeclareSectionCommand[
    beforeskip=0pt
]{chapter}
\clearpairofpagestyles
\ohead{\headmark}
\ofoot{\pagemark}

% Seitenzaehlung wie im Vorspann von master.pdf
\pagenumbering{roman}
\setcounter{page}{6}

\begin{document}

\onehalfspacing

\chapter*{Kurzfassung}
% In master.pdf schrumpft der Abstand nach der Kapitelueberschrift durch
% \flushbottom um 3,0605 pt. Hier ausgeglichen, damit der Textblock exakt
% an derselben Stelle steht wie in master.pdf.
\vspace*{-3.072pt}
Die Stichprobenanalyse im Internal Audit erfordert die manuelle Prüfung zahlreicher Rechnungen und Freigaben
und ist entsprechend zeitaufwendig. Large Language Models können diese dokumentenbasierte Prüfung unterstützen.
Ihre Neigung zu Halluzinationen und nicht deterministischen Ausgaben wirft im Audit jedoch Fragen hinsichtlich des zuverlässigen Einsatzes auf.
Die vorliegende Arbeit untersucht daher, wie sich verschiedene Large Language Models hinsichtlich
ihrer Eignung zur Unterstützung von Stichprobenanalysen im Audit-Bereich der Freudenberg Group unterscheiden.

Hierzu wurden die acht in Microsoft Copilot Studio verfügbaren Modelle von OpenAI (GPT-4.1, GPT-5, GPT-5.3, GPT-5.5)
und Anthropic (Claude Sonnet 4.6, Opus 4.6, Opus 4.7, Opus 4.8) jeweils als Agent mit identischen,
an den Prüfschritten des Corporate Audit ausgerichteten Instruktionen konfiguriert. Als Testdatensatz dienen
100 künstlich erzeugte Rechnungen mit zugehörigen Freigabenachweisen. In diese wurden 24 Audit Exceptions aus
sechs Fehlerkategorien eingebaut. Bewertet wurden die Modelle über eine Nutzwertanalyse mit den Kriterien
Recall (40\,\%), Precision (30\,\%) und Konsistenz (30\,\%). Die Konsistenz wurde anhand von fünf Durchläufen
über 30 Rechnungen gemessen und die Robustheit der Gewichtung über eine Sensitivitätsanalyse geprüft.

Den höchsten Nutzwert erreicht Claude Opus 4.7 mit 92,9 von 100 Punkten (Recall 91,7\,\%, Precision 95,7\,\%),
gefolgt von Opus 4.8 mit 89,6 Punkten. Nur diese beiden Modelle erzielen bei Recall, Precision und Konsistenz
durchgehend Werte über 80\,\%. Die Modelle von Anthropic belegen die Plätze eins bis vier und liegen bei Recall
und Konsistenz vor allen Modellen von OpenAI. Diese weisen dagegen jeweils eine deutliche Schwäche auf. GPT-4.1 übersieht
fast die Hälfte der Exceptions, während bei GPT-5, GPT-5.3 und GPT-5.5 mehr als die Hälfte der Meldungen unberechtigt ist.
Die Rangfolge bleibt in allen Szenarien der Sensitivitätsanalyse stabil. Neuere Modellversionen schneiden dabei
nicht automatisch besser ab als ältere.

Für die Freudenberg Group wird der Einsatz von Opus 4.7 oder Opus 4.8 als Agent zur Vorprüfung empfohlen.
Da auch das beste Modell nicht alle Exceptions erkennt, dienen die Ergebnisse als Entscheidungshilfe,
und die abschließende Bewertung verbleibt beim Auditor.

\end{document}
